% Cote 119: offprints, reports, courses, CV and letters (typed and printed part): archive cote 119, site page 266 (PDF page 267).
% typescript. Draft compared: none.
% AI-transcribed by Claude Opus 5.5 (claude-opus-5-5, adaptive thinking, effort high), Message Batches API, batch r119_full (msgbatch_01DwZNm5Y7YsKBWyAN4sFs8w), 2026-11-13; not yet checked by eye.
\scanpage{119/266}
\mspage{13}

ou pédagogiques, etc. Ces mêmes éléments d'appréciation auront évidemment leur
rôle à jouer pour la nomination d'un docteur d'enseignement à un poste de
maître de conférences (mention enseignement), décidée (rappelons-le) par
l'Assemblée de département, donc avec participation effective des délégués
d'étudiants. (Pour la nomination d'un maître de conférences mention recherche,
il conviendrait par contre de prévoir, en cas de demande par l'Institut de
Mathématique, une cooptation par une assemblée plus restreinte, comprenant
tous les membres de l'Institut de Mathématique, et où ne figuereraient pas de
délégués étudiants).

5. Le point le plus délicat à l'heure actuelle semble celui du \uline{rôle dévolu à
l'assistant et au maître-assistant} dans l'Université de demain. La doctrine
la plus satisfaisante dans l'absolu semblerait la suivante. \uline{La fonction d'assis-%
tant et de maître-assistant serait considéré comme une fonction essentiellement}
non permanente, \uline{correspondante aux dernières étapes de la formation d'un
maître de conférences} (mention enseignement, le plus souvent). Elle doit donc
\uline{normalement} déboucher sur un doctorat (d'enseignement, le plus souvent), et
sur une nomination de maître de conférences, après une période de l'ordre de
quatre à six ans. Dans cette optique, il convient donc de ne titulariser un
assistant stagiaire que lorsque le département a de sérieuses raisons de penser
que le candidat a les qualités requises pour aller jusqu'au doctorat, donc
après la soutenance de la thèse, et sans doute pas avant deux ans de pratique
de l'enseignement le plus souvent. Il faudrait donc pourvoir à la réorientation
professionnelle des stagiaires assistants qui ne donneraient pas satisfaction,
et, exceptionnellement, des enseignants ayant exercé les fonctions d'assistant
ou maître-assistant pendant six années sans arriver à obtenir le doctorat. Ce
dernier cas sera évidemment délicat du fait du statut de la fonction publique
dont bénéficiera cet enseignant, mais on peut espérer qu'il sera rare. et
qu'il sera possible dans la plupart des cas, avec l'accord de l'intéressé, de
lui trouver un débouché en-dehors de l'enseignement, plus en rapport avec ses
aptitudes réelles.

Si on compte que la durée moyenne des fonctions d'assistants ou
maître-assistant avant le doctorat et la nomination à un poste de maître de
conférences est de cinq années, sur une activité d'enseignement qui porte
sur environ quarante ans du total (de vingt-cinq à soixante-cinq ans), on

\margin{266}
